\chapter{خريطة الجبهات البحثية الحديثة}

\section{وظيفة هذا الجزء}

هذا الجزء ليس قائمة أخبار، بل مراجعة دورية للحدود التي تتقاطع عندها تقنيات نظرية الأعداد التحليلية مع مسائل لم تُحسم. كل قسم حديث سيحدد: المسألة، وأفضل نتيجة معروفة في تاريخ النسخة، والأداة الحاسمة، والعوائق، والمراجع الأصلية.

\section{محاور أولية}

\begin{enumerate}
  \item أصفار وقيم دوال زيتا ودوال \(L\)، بما في ذلك العزوم والقيم المتطرفة.
  \item الفجوات الصغيرة والكبيرة بين الأعداد الأولية.
  \item التوزيع في المتتاليات الحسابية والفترات القصيرة.
  \item مسائل دون التحدب وقيم دوال \(L\) في النقطة المركزية.
  \item صيغ التتبع النسبية والتبادلية الطيفية.
  \item الدوال الضربية الادعائية والارتباطات من نمط تشاولا.
  \item فك الاقتران وتطبيقاته في المجاميع الأسية ومسألة وارينغ.
  \item الإحصاء الحسابي وعد الحقول والفئات المثالية.
  \item المصفوفات العشوائية ونماذج إحصاء الأصفار والقيم.
  \item الحساب الموثق عالي الدقة والتحقق العددي من التخمينات.
\end{enumerate}

\section{بروتوكول التحديث}

لكل تحديث بحثي يسجل تاريخ القطع، ومصدر أولي، وحالة النتيجة: مثبتة، أو مشروطة، أو معلنة قبل التحكيم، أو تخمينًا. ويجب تجنب وصف أي نتيجة بأنها ``الأحدث'' دون بحث محدث في تاريخ الإصدار.
